\chapter{Theoretischer Hintergrund und Stand der Forschung}

Frühe Ansätze der Insolvenzprognose beruhen auf Kennzahlen und Linearkombinationen. 
Beavers univariate Tests identifizierten bereits in den 1960er-Jahren Trennschärfe in Liquiditäts- und Ertragskennzahlen, während der Z-Score nach \textcite{altman1968financial} eine multivariate 
Diskriminanzanalyse etablierte. Diese Klassiker wirken bis heute fort, stoßen aber bei 
Heterogenität, Nichtlinearität und zeitlicher Verschiebung der Signale an Grenzen.

Mit der Ausbreitung datengetriebener Methoden erweiterte sich das Spektrum auf 
Logit-Modelle, Boosting/Bagging-Ansätze und tiefe Netze. Meta-Analysen und 
Vergleichsarbeiten zeigen, dass moderne ML-Methoden—insbesondere Ensemble-Verfahren—
klassische Spezifikationen häufig übertreffen, wenn Vorverarbeitung, Regularisierung 
und Validierung sorgfältig ausgelegt sind \parencite{barboza2017machine,sun2024bankruptcy,hastieelementsstatistical2009}.

Die ökonomische Intuition bleibt jedoch zentral: Kennzahlen zu Profitabilität, Liquidität, 
Leverage, Effizienz und Aktivität bilden die klassische Grundlage der Früherkennung. 
Gleichzeitig erfordern Verhältniskennzahlen besondere Sorgfalt in der Aufbereitung 
(Winsorisierung/Imputation) und in der Modellierung (Robustheit gegenüber Ausreißern 
und Verzerrungen) \parencite{vonhippel2013passive,vanbuuren2018fimd}. Multikollinearität ist 
weit verbreitet und verzerrt Koeffizienten und Unsicherheiten, weshalb VIF-Grenzen als 
praktikable Daumenregel genutzt werden—unter Beachtung der methodischen Einschränkungen \parencite{obrien2007caution,pennstate2024vif}.

Methodisch ist die Validierung der entscheidende Hebel zur Vermeidung von Optimismus. 
Nested Cross-Validation reduziert Selektionsbias, der bei Hyperparameter- und 
Merkmalsauswahl sonst systematisch entsteht \parencite{cawley2010overfitting,varma2006bias}. 
Für unbalancierte Zielgrößen empfiehlt sich neben ROC-AUC die PR-AUC \parencite{saito2015precision}. 
Zudem quantifiziert ein Stabilitätsmaß die Robustheit der Merkmalsauswahl über Folds 
\parencite{nogueira2018stability}, und die 
EPV-Regel dient als Guardrail gegen überparametrisierte Logit-Modelle \parencite{peduzzi1996epv,steyerberg2019clinical}. 
Schließlich mahnen Arbeiten zur Interpretierbarkeit, für risikorelevante Entscheidungen 
nach Möglichkeit transparente Modelle oder stabile, kleine Feature-Sets vorzuziehen \parencite{rudin2019stop}.

Vor diesem Hintergrund verbindet die vorliegende Arbeit klassische Bilanzlogik mit 
aktuellen ML-Praktiken: kontrollierte Multikollinearitätsreduktion, 
mehrgleisige Feature-Selektion mit Stabilität und Nested-CV, EPV-konforme 
Endmengen sowie ein fairer, stratifizierter Modellvergleich über mehrere 
Prognosehorizonte.
